
% ============================================================
% SECTION 03
% INVARIANCE AND DIFFERENCE
% Chapter 02: Invariance
% Book I: Ontology of Viability
% ============================================================

\section{Invariance and Difference}
\label{sec:invariance-and-difference}

% ------------------------------------------------------------
% 3.1. THE PROBLEM OF DIFFERENCE IN THE ANALYSIS
% OF INVARIANCE
% ------------------------------------------------------------

The concept of invariance cannot be adequately examined
without considering its relationship to difference.
An invariance claim identifies something that remains
unchanged in a particular respect within a specified
domain of changes or transformations. The specification
of such a domain necessarily raises the question of
what changes, what differs, and how the relevant
differences are distinguished from what remains
unchanged.

Difference is therefore indispensable to the analytical
formulation of invariance. However, this does not mean
that difference is ontologically prior to invariance,
or that invariance is merely an effect produced by
the comparison of different states.

The relationship must be approached without assuming
that every difference negates invariance, or that
every instance of invariance eliminates difference.
These assumptions would reduce two distinct concepts
to mutually exclusive conditions.

An existent may exhibit differences in certain
determinations while remaining invariant in another
respect. The task is to clarify the conditions under
which these two descriptions can be applied without
contradiction.

The present section therefore examines the different
ways in which difference may be related to invariance
and establishes the conceptual restrictions necessary
for a coherent account of both.

% ------------------------------------------------------------
% 3.2. DIFFERENCE AS A SPECIFIED DETERMINATION
% ------------------------------------------------------------

Difference must first be distinguished from the
undifferentiated notion of change.

To say that two conditions differ is to identify
a distinction between them with respect to some
specified determination. The determination may
concern position, quantity, configuration, relation,
quality, or another aspect relevant to the analysis.

Difference is therefore not adequately expressed
by the mere assertion that two conditions are not
identical. A meaningful analysis must clarify
the respect in which the difference is being
considered.

For example, an existent may occupy different
spatial positions under two conditions. The
positions differ, but this difference does not
necessarily establish that every other determination
of the existent also differs.

Likewise, two configurations may differ in their
arrangement while preserving a specified relation
among their components.

These examples indicate that difference is often
determination-specific. A difference concerning
one aspect does not automatically extend to all
other aspects.

This restriction is essential for the analysis
of invariance because a claim of invariance must
identify the particular determination that remains
unchanged despite the differences occurring
within the specified domain.

% ------------------------------------------------------------
% 3.3. DIFFERENCE DOES NOT NECESSARILY NEGATE
% INVARIANCE
% ------------------------------------------------------------

The existence of difference does not necessarily
exclude invariance.

If an existent undergoes a transformation in which
one determination changes while another remains
unchanged, the two descriptions refer to different
respects of the same case.

The difference and the invariance are not contradictory
because they do not concern the same determination
under the same conditions.

This can be expressed schematically. Let \(A\)
represent an existent, and let \(T\) represent
a specified transformation. Let \(F\) and \(G\)
represent two distinct determinations considered
in relation to that existent.

It is possible that:

\[
F(T(A)) \neq F(A)
\]

while:

\[
G(T(A)) = G(A)
\]

The first expression indicates variation in
determination \(F\), while the second indicates
invariance in determination \(G\), within the
specified transformation.

These expressions are schematic and presuppose
that the relevant determinations and transformation
are sufficiently defined. They do not establish
that every ontological difference can be represented
mathematically in this way.

Their purpose is to clarify that difference and
invariance may coexist when they concern distinct
respects of an existent.

Consequently, the presence of difference cannot
be treated as sufficient evidence against every
claim of invariance.

% ------------------------------------------------------------
% 3.4. DIFFERENCE WITHIN THE SAME DETERMINATION
% ------------------------------------------------------------

A more difficult problem arises when the same
determination is considered under different
conditions.

If a particular determination changes under
a specified transformation, it cannot simultaneously
be described as unchanged under that same
transformation and in the same respect.

The relevant distinction is between two different
claims:

\begin{enumerate}
    \item A determination changes in relation
    to a specified transformation.

    \item A determination remains unchanged
    in relation to that same transformation.
\end{enumerate}

If both claims refer to precisely the same
determination, transformation, and respect,
they are incompatible.

However, the incompatibility of these claims
does not establish that the existent as a whole
has changed in every respect or that its identity
has been lost.

It only establishes that the particular
determination under consideration is not
invariant within the specified domain.

The analysis must therefore avoid transferring
a conclusion about one determination to the
entire existent without an additional argument.

% ------------------------------------------------------------
% 3.5. DIFFERENCE BETWEEN EXISTENTS AND
% DIFFERENCE WITHIN AN EXISTENT
% ------------------------------------------------------------

The concept of difference operates in more than
one analytical context.

First, difference may concern the distinction
between two or more existents. In this case,
the analysis concerns how one existent differs
from another.

Second, difference may concern distinct conditions
or determinations associated with one existent.
In this case, the analysis concerns variation
within the consideration of a single existent.

These two contexts must not be confused.

The distinction between two existents is not
equivalent to the difference between two conditions
of one existent. Likewise, the variation of
a determination within one existent does not
automatically establish a difference between
two independent existents.

The relational definition of invariance must
therefore specify whether its domain concerns
transformations of one existent, comparisons
between different existents, or another form
of variation.

Without this specification, an invariance claim
may confuse the persistence of a determination
within one existent with the similarity or
difference between multiple existents.

The ontological status of relations between
different existents, and the extent to which
such relations contribute to their identity,
remain questions requiring further examination.

% ------------------------------------------------------------
% 3.6. DIFFERENCE, COMPARISON, AND THE
% POSITION OF THE KNOWER
% ------------------------------------------------------------

Difference is frequently recognized through
comparison. A knower considers two or more
conditions and identifies a determination
that differs between them.

This comparative activity is epistemic.
It depends upon the distinctions available
to the knower, the criteria of comparison,
and the domain within which the comparison
is conducted.

However, the epistemic recognition of a
difference must not be confused with the
ontological existence of the difference itself.

A difference may exist without being recognized
by a particular knower. Conversely, an apparent
difference may result from an inadequate
description, an inappropriate comparison,
or a mistaken identification of the relevant
determinations.

The same restriction applies to invariance.
The recognition of something as unchanged
depends upon the identification of the relevant
domain and respect of comparison. Yet this
epistemic dependence does not establish that
invariance is created by the act of recognition.

Knowledge and non-knowledge are intertwined
within the situated condition of the knower.
The recognition of one difference may coexist
with ignorance of other differences, just as
the identification of one invariant determination
may coexist with the inability to recognize
other invariant or variable aspects.

This limitation must not be interpreted as
a deficiency that can be eliminated simply
by assuming an unrestricted observer. The
knower is itself situated within being and
does not occupy a position outside the
ontological conditions being investigated.

% ------------------------------------------------------------
% 3.7. DIFFERENCE AND THE SCOPE OF AN
% INVARIANCE CLAIM
% ------------------------------------------------------------

An invariance claim is always restricted by
the domain within which it is formulated.

A determination may remain unchanged under
one set of transformations while varying
under another. Consequently, the identification
of invariance within a limited domain cannot
be generalized automatically to every possible
domain of change.

The relationship can be represented schematically.

Let \(D_1\) and \(D_2\) represent two distinct
domains of transformations. A determination
\(F\) may satisfy:

\[
F(T(A)) = F(A)
\qquad
\text{for } T \in D_1
\]

without satisfying the same condition for
transformations belonging to \(D_2\).

The expression does not establish that the
two domains are necessarily independent or
that every transformation can be assigned
to a clearly bounded mathematical set. It
illustrates the limited scope of an invariance
claim.

The same determination may therefore be
invariant relative to one domain and variable
relative to another.

This possibility prevents the concept of
invariance from being interpreted as an
unrestricted assertion of permanence.

It also establishes the need to identify
the domain of comparison before making
general claims about what remains unchanged.

% ------------------------------------------------------------
% 3.8. DIFFERENCE DOES NOT NECESSARILY
% IMPLY THE LOSS OF IDENTITY
% ------------------------------------------------------------

The distinction between difference and
invariance has direct implications for
the analysis of identity.

If an existent changes in one or more
determinations, it does not follow automatically
that its identity has ceased to be.

The previous chapter established that
existence and identity are inseparable.
The present analysis must preserve this
commitment without prematurely explaining
how identity relates to every possible
transformation.

A difference between two conditions of
an existent may be compatible with the
continuation of its identity. However,
the existence of such compatibility in
some cases does not establish a universal
rule governing every transformation.

The precise relationship between identity,
difference, and the persistence of an
existent remains dependent upon questions
that cannot be settled by the concept of
invariance alone.

In particular, the present section does
not determine which changes are compatible
with the persistence of identity, what
conditions may constitute a loss of identity,
or whether these questions have the same
answer across different ontological domains.

These issues must remain open until the
broader structure of the ontology has
been sufficiently developed.

% ------------------------------------------------------------
% 3.9. DIFFERENCE AS A CONDITION OF
% CONCEPTUAL DISTINCTION
% ------------------------------------------------------------

The analysis of invariance requires the
ability to distinguish what varies from
what remains unchanged. Difference therefore
has an indispensable analytical function.

Without identifying relevant differences,
it would be impossible to specify the
domain within which invariance is being
asserted.

However, this analytical function must
not be confused with an ontological claim
that difference produces invariance.

The identification of a distinction makes
an invariance claim conceptually articulable.
It does not necessarily establish the
ontological origin of the invariant
determination.

Likewise, the fact that a determination
is identified through comparison does
not mean that its existence depends upon
the comparison.

The analysis must consequently preserve
three distinctions:

\begin{enumerate}
    \item Difference as an ontological
    distinction or variation.

    \item Comparison as an epistemic
    operation through which differences
    may be recognized.

    \item Invariance as a relation in which
    a specified determination remains
    unchanged within a defined domain
    of change.
\end{enumerate}

These concepts are related, but none
can be reduced to another without
additional argument.

% ------------------------------------------------------------
% 3.10. SYNTHESIS AND TRANSITION
% ------------------------------------------------------------

The relationship between invariance and
difference can now be stated more precisely.

Difference identifies variation or
distinction with respect to a specified
determination. Invariance identifies
the persistence of a particular
determination in relation to a specified
domain of changes or transformations.

Difference does not necessarily negate
invariance because an existent may vary
in one respect while remaining unchanged
in another.

At the same time, the persistence of
one determination does not eliminate
the differences occurring in other
determinations.

The analysis therefore rejects two
reductions:

\begin{enumerate}
    \item The reduction of invariance
    to the absence of all difference.

    \item The reduction of difference
    to the necessary negation of
    every form of invariance.
\end{enumerate}

The relational definition adopted in
this chapter makes it possible to
examine both concepts within a specified
domain without treating them as
mutually exclusive.

Nevertheless, the present analysis
does not explain the ontological origin
of either difference or invariance.
Nor does it establish whether every
existent necessarily exhibits invariant
determinations.

These questions remain open for the
subsequent development of the ontology
of viability.

The next section examines the relation
between invariance and change. Although
difference and change are closely
connected, they are not identical.
A more precise analysis of their
relationship is necessary to determine
how invariance can be understood
in relation to transformation.

% ============================================================
% END OF SECTION 03
% ============================================================
